\documentclass[a4paper,11pt]{ltjsarticle}

\usepackage{luatexja}
\usepackage{geometry}
\geometry{margin=25mm}

\usepackage{array}
\usepackage{longtable}
\usepackage[table]{xcolor}
\usepackage{booktabs}

\usepackage{enumitem}

\usepackage{hyperref}

\usepackage{titlesec}
\titleformat{\section}{\large\bfseries}{}{0em}{}

\usepackage{listings}
\lstset{
  basicstyle=\ttfamily\small,
  frame=single,
  breaklines=true,
  numbers=left,
  numberstyle=\tiny\color{gray},
  keepspaces=true,
  showstringspaces=false,
  columns=fullflexible,
  tabsize=2,
}

\usepackage{tikz}
\usetikzlibrary{shapes.geometric, arrows.meta, positioning, calc, fit}

\usepackage{float}

\definecolor{ganttCell}{HTML}{4FC3D0}
\definecolor{ganttEarly}{HTML}{FFB74D}

\providecommand{\％}{\%}

\setlength{\parindent}{0pt}
\setlist[itemize]{leftmargin=2em}

\begin{document}

\begin{center}
    {\LARGE \textbf{作業ログ（第7週末）}}\\[1em]
\end{center}

\begin{tabular}{|p{4cm}|p{10cm}|}
\hline
報告者 & 櫛濱 和輝 (25G1046)\\
\hline
報告日 & 2026年6月9日\\
\hline
プロジェクト名 & 赤外線通信で歌うカエル\\
\hline
担当パート & Processing（ピアノ演奏プログラム）／Arduino（ドラムのプログラム \texttt{hack-drum.ino}）\\
\hline
本来の今週の位置づけ & 第7週末（6/9）．ガントチャート上，第7週（6/3〜6/9）に櫛濱の割当タスクはない\\
\hline
報告対象 & 第6週末（6/2）までに完了した「ピアノのプログラム（Processing）」と「ドラムのプログラム \texttt{hack-drum.ino}（GC の 240）」．ドラム音色 170/171 の \texttt{.pde} は久家担当の別タスク\\
\hline
補足 & 第6週はログ未提出，かつ GC 上で第7週は櫛濱の割当タスクがないため，本（第7週末）ログで第6週末までの成果をまとめて報告\\
\hline
\end{tabular}

\vspace{1em}

%=====================================================
\section{1. 進捗サマリー}

\begin{itemize}
    \item 現在地：第7週末（6/9）．ガントチャート（§2.2）上，第7週（6/3〜6/9）に櫛濱の割当タスクはないため，本ログでは\textbf{第6週末（6/2）までに完了した作業}を報告する
    \item 第5〜6週で \textbf{ピアノのプログラム（Processing：110 音色＋120 受信＋210 子機連携）} と \textbf{ドラムのプログラム（\texttt{hack-drum.ino}，GC の 240）} を実装．ピアノは\textbf{実機で動作確認まで完了}，ドラムは\textbf{作成のみで実機テストは未実施}
    \item 第6週は対面できなかったため検証範囲が限られるが，ピアノ系（親機 → 赤外線 → ピアノ子機 → Processing）では \texttt{hack-oya2.ino} と \texttt{hack\_ko2.pde} を用いて\textbf{親機からの譜面送信と発音に成功}（第6週末までに確認）
    \item ドラムは久家と作業を分け，櫛濱が \texttt{hack-drum.ino}（GC の 240：ドラムのプログラム．NEC 受信→シリアル中継）を作成（ドラム音色 170/171 の \texttt{.pde} は久家担当の別タスク）．本プログラムは\textbf{作成したのみで実機テストは未実施}
    \item 全体進捗率：45\％（音色作成フェーズ完了，子機・通信フェーズ進行中）
    \item 主要成果：
    \begin{itemize}
        \item Processing 側（\texttt{hack\_ko2.pde}）が子機 Arduino からシリアルで届く MIDI 番号を受信し，\texttt{PianoInstrument} で発音．MIDI→周波数変換 $f=440\cdot2^{(m-69)/12}$ を実装
        \item ピアノ系では「親機 → 赤外線（NEC）→ ピアノ子機 → USB シリアル → Processing → 発音」の経路が一気通貫で動作し，譜面の送信を実機で確認
        \item ドラムのプログラム（\texttt{hack-drum.ino}，GC の 240）は NEC フレーム（\texttt{address==2}）のみを拾い \texttt{command}（MIDI 番号）をシリアルへ中継する設計で実装（自前タイマを持たず\textbf{テンポは親機任せ}）．久家のドラム音色（\texttt{.pde}）と組み合わせて鳴らす想定だが，\textbf{実機テストは未実施}
    \end{itemize}
    \item 重大課題（影響度：高）：
    \begin{itemize}
        \item ドラムのプログラム（\texttt{hack-drum.ino}，GC の 240）が未テスト：第6週に対面できず，久家のドラム音色（\texttt{.pde}）と結合した実機での受信・発音確認ができていない
        \item 音切り替え時のノイズ：現状 \texttt{noteOff()} が \texttt{unpatch} の即時切断のため，MOP5（フェードアウト）が未対応
        \item 赤外線通信の成功率（MOP3：95\％以上）が未計測．取りこぼし時の挙動も未検証
    \end{itemize}
\end{itemize}

%=====================================================
\section{2. ガントチャート}

計画書 v1.2（2026/5/29 版）でチーム全体のスケジュールが更新されたため，チーム全体・個人の両方について変更前後を併記する．着色セル（\colorbox{ganttCell}{\phantom{xx}}）は当該タスクの実施週を表す．

\renewcommand{\arraystretch}{1.3}

\begingroup
\setlength{\tabcolsep}{3pt}

\subsection*{2.1 チーム全体ガントチャート（変更前：第5週末提出版）}

\begin{longtable}{|p{0.9cm}|p{2.0cm}|p{4.0cm}|p{1.5cm}|c|c|c|c|c|}
\hline
\rowcolor[gray]{0.9}
管理 & 第2層 & 第3層タスク & 担当者 & 6週 & 7週 & 8週 & 9週 & 10週 \\
\hline
110 & 100 (Processing) & ピアノ音の作成 & 櫛濱和輝 & \cellcolor{ganttCell} & & & & \\
\hline
120 & 100 (Processing) & arduinoとの通信用関数（ピアノ） & 櫛濱和輝 & & & \cellcolor{ganttCell} & \cellcolor{ganttCell} & \\
\hline
130 & 100 (Processing) & フルート音の作成 & 家田祐希 & \cellcolor{ganttCell} & & & & \\
\hline
140 & 100 (Processing) & arduinoとの通信用関数（フルート） & 家田祐希 & & & \cellcolor{ganttCell} & \cellcolor{ganttCell} & \\
\hline
150 & 100 (Processing) & トランペット音の作成 & 木下遼介 & \cellcolor{ganttCell} & & & & \\
\hline
160 & 100 (Processing) & arduinoとの通信用関数（トランペット） & 木下遼介 & & & \cellcolor{ganttCell} & \cellcolor{ganttCell} & \\
\hline
170 & 100 (Processing) & ドラム音の作成（キック） & 久家力孔 & \cellcolor{ganttCell} & & & & \\
\hline
171 & 100 (Processing) & ドラム音の作成（シンバル） & 久家力孔 & \cellcolor{ganttCell} & & & & \\
\hline
180 & 100 (Processing) & arduinoとの通信用関数（ドラム） & 久家力孔 & & & \cellcolor{ganttCell} & \cellcolor{ganttCell} & \\
\hline
210 & 200 (子機) & ピアノのプログラム & 櫛濱和輝 & & \cellcolor{ganttCell} & & & \\
\hline
220 & 200 (子機) & フルートのプログラム & 家田祐希 & & \cellcolor{ganttCell} & & & \\
\hline
230 & 200 (子機) & トランペットのプログラム & 木下遼介 & & \cellcolor{ganttCell} & & & \\
\hline
240 & 200 (子機) & 楽譜の設定 & 渡邊乃斗 & \cellcolor{ganttCell} & & & & \\
\hline
310 & 300 (親機) & 赤外線管理プログラム & 渡邊乃斗 & & \cellcolor{ganttCell} & & & \\
\hline
320 & 300 (親機) & テンポ管理プログラム & 渡邊乃斗 & & & \cellcolor{ganttCell} & \cellcolor{ganttCell} & \\
\hline
410 & 400 (その他) & 回路作成 & 手代木巧 & \cellcolor{ganttCell} & & & & \\
\hline
420 & 400 (その他) & LEDマトリクス作成 & 手代木巧 & & \cellcolor{ganttCell} & \cellcolor{ganttCell} & \cellcolor{ganttCell} & \\
\hline
510 & 500 (包括) & 結合テスト & 全員 & & & & & \cellcolor{ganttCell} \\
\hline
\end{longtable}

\subsection*{2.2 チーム全体ガントチャート（変更後：計画書 v1.2）}

計画書 v1.2（2026/5/29 版）のチーム全体ガントチャートに合わせて更新した．担当者は楽器グループごとの 2 名体制に整理し，子機側プログラムを 210〜260 として再構成している．

\begin{longtable}{|p{6.2cm}|p{2.3cm}|c|c|c|c|c|c|}
\hline
\rowcolor[gray]{0.9}
第3層：活動タスク & 担当者 & 5週 & 6週 & 7週 & 8週 & 9週 & 10週 \\
\hline
110:ピアノ音の作成 & 櫛濱・久家 & \cellcolor{ganttCell} & \cellcolor{ganttCell} & & & & \\
\hline
120:arduinoとの通信用関数（ピアノ） & 櫛濱・久家 & \cellcolor{ganttCell} & \cellcolor{ganttCell} & & & & \\
\hline
130:フルート音の作成 & 家田・渡邊 & \cellcolor{ganttCell} & \cellcolor{ganttCell} & & & & \\
\hline
140:arduinoとの通信用関数（フルート） & 家田・渡邊 & \cellcolor{ganttCell} & \cellcolor{ganttCell} & & & & \\
\hline
150:トランペット音の作成 & 木下・手代木 & \cellcolor{ganttCell} & \cellcolor{ganttCell} & & & & \\
\hline
160:arduinoとの通信用関数（トランペット） & 木下・手代木 & \cellcolor{ganttCell} & \cellcolor{ganttCell} & & & & \\
\hline
170:ドラム音の作成（キック） & 櫛濱・久家 & & \cellcolor{ganttCell} & & & & \\
\hline
171:ドラム音の作成（シンバル）\textsuperscript{新規} & 櫛濱・久家 & & \cellcolor{ganttCell} & & & & \\
\hline
180:arduinoとの通信用関数（ドラム） & 櫛濱・久家 & & \cellcolor{ganttCell} & & & & \\
\hline
210:ピアノのプログラム & 櫛濱・久家 & \cellcolor{ganttCell} & & & & & \\
\hline
220:フルートのプログラム & 家田・渡邊 & \cellcolor{ganttCell} & & & & & \\
\hline
230:トランペットのプログラム & 木下・手代木 & \cellcolor{ganttCell} & & & & & \\
\hline
240:ドラムのプログラム & 櫛濱・久家 & & \cellcolor{ganttCell} & & & & \\
\hline
260:楽譜進行プログラム & 家田・渡邊 & & \cellcolor{ganttCell} & \cellcolor{ganttCell} & & & \\
\hline
310:赤外線管理プログラム & 家田・渡邊 & & & \cellcolor{ganttCell} & & & \\
\hline
320:テンポ管理プログラム & 家田・渡邊 & & & \cellcolor{ganttCell} & \cellcolor{ganttCell} & & \\
\hline
410:回路作成 & 木下・手代木 & & \cellcolor{ganttCell} & & & & \\
\hline
420:LEDマトリクス作成 & 木下・手代木 & & \cellcolor{ganttCell} & \cellcolor{ganttCell} & & & \\
\hline
510:結合テスト & 全員 & & & & & \cellcolor{ganttCell} & \cellcolor{ganttCell} \\
\hline
\end{longtable}

\endgroup

\textbf{主な変更点（第5週末提出版 → 計画書 v1.2）：}
\begin{itemize}
    \item 担当を楽器グループごとの 2 名体制（櫛濱・久家／家田・渡邊／木下・手代木）に整理
    \item 音色作成（110〜160）と通信用関数（120/140/160）を第5〜6週に集約
    \item 子機プログラム（210/220/230）を第7週 → 第5週へ前倒し
    \item ドラム子機を 240：ドラムのプログラム（第6週）として新設し，旧 240（楽譜の設定）を 260：楽譜進行プログラム（第7週）へ整理
    \item 171（ドラムシンバル）は新規追加タスク
    \item 結合テスト（510）を第9〜10週に設定
\end{itemize}

\subsection*{2.3 個人ガントチャート（変更前）}

第5週末ログ時点の個人計画．210（ピアノのプログラム）は第7週に割り当てられていた．

\begin{longtable}{|p{0.9cm}|p{2.0cm}|p{6.0cm}|c|c|c|c|c|c|}
\hline
\rowcolor[gray]{0.9}
管理 & 第2層 & 第3層タスク & 5週 & 6週 & 7週 & 8週 & 9週 & 10週 \\
\hline
110 & 100 (Processing) & ピアノ音の作成 & \cellcolor{ganttCell} & \cellcolor{ganttCell} & & & & \\
\hline
210 & 200 (子機) & ピアノのプログラム & & & \cellcolor{ganttCell} & & & \\
\hline
120 & 100 (Processing) & arduinoとの通信用関数（ピアノ） & & & & \cellcolor{ganttCell} & \cellcolor{ganttCell} & \\
\hline
510 & 500 (包括) & 結合テスト（全員） & & & & & & \cellcolor{ganttCell} \\
\hline
\end{longtable}

\subsection*{2.4 個人ガントチャート（変更後）}

計画書 v1.2 のチーム全体ガント（§2.2）から櫛濱担当分を抜き出した個人計画．音色作成（110）・通信用関数（120）を第5〜6週，ピアノのプログラム（210）を第5週，ドラムのプログラム（240，\texttt{hack-drum.ino}）を第6週に置く．なおドラムは第6週に対面できず\textbf{実機テストは未実施}（プログラム作成のみ）で，音色 170/171 の \texttt{.pde} は久家担当の別タスク．

\begin{longtable}{|p{0.9cm}|p{2.0cm}|p{6.0cm}|c|c|c|c|c|c|}
\hline
\rowcolor[gray]{0.9}
管理 & 第2層 & 第3層タスク & 5週 & 6週 & 7週 & 8週 & 9週 & 10週 \\
\hline
110 & 100 (Processing) & ピアノ音の作成 & \cellcolor{ganttCell} & \cellcolor{ganttCell} & & & & \\
\hline
120 & 100 (Processing) & arduinoとの通信用関数（ピアノ） & \cellcolor{ganttCell} & \cellcolor{ganttCell} & & & & \\
\hline
210 & 200 (子機) & ピアノのプログラム & \cellcolor{ganttCell} & & & & & \\
\hline
240 & 200 (子機) & ドラムのプログラム（\texttt{hack-drum.ino}） & & \cellcolor{ganttCell} & & & & \\
\hline
510 & 500 (包括) & 結合テスト（全員） & & & & & \cellcolor{ganttCell} & \cellcolor{ganttCell} \\
\hline
\end{longtable}

\subsection*{2.5 今週の作業位置づけ}

\begin{itemize}
    \item 現在地は第7週末（6/9）．GC 上，第7週（6/3〜6/9）に櫛濱の割当タスクはないため，本ログは第6週末（6/2）までの成果を報告する．
    \item ピアノのプログラム（110 音色＋120 受信＋210 子機連携）は第5〜6週で実装し，親機 → 赤外線 → ピアノ子機 → Processing の経路で\textbf{譜面送信・発音に成功}．
    \item ドラムは久家と作業を分け，櫛濱が 240：ドラムのプログラム（\texttt{hack-drum.ino}）を第6週に作成したが，第6週に対面できず\textbf{実機テストは未実施}（音色 170/171 は久家担当）．
    \item 第8週（6/10〜）以降は，ドラムの実機テスト，120（通信用関数・ピアノ）の正式化，MOP5（フェードアウト）対応に移る．
\end{itemize}

%=====================================================
\section{3. 作業ログ}

\renewcommand{\arraystretch}{1.4}

\begin{longtable}{|p{2.6cm}|p{1.4cm}|p{1.6cm}|p{1.4cm}|p{1cm}|p{2.4cm}|p{3.0cm}|}
\hline
\rowcolor[gray]{0.9}
作業項目 & 開始日 & 予定完了日 & 状態 & 進捗率 & 依存関係・ブロッカー & コメント \\
\hline
ピアノのプログラム (210/110/120) & 5/22 & 5/27 & 完了 & 100\％ & 子機 Arduino（\texttt{hack-ko2.ino}）／親機 & Processing でシリアル受信した MIDI 番号を周波数へ変換し \texttt{PianoInstrument} で発音．波形描画も実装（付録 9.3.1） \\
\hline
ドラムのプログラム (240) \texttt{hack-drum.ino} & 5/28 & 6/2 & 作成済（未検証） & 実装100\％／検証0\％ & 親機の NEC 送出／久家のドラム音色 \texttt{.pde} & NEC（\texttt{address==2}）の \texttt{command} をシリアルへ中継する設計で作成．第6週に対面できず実機テスト未実施（付録 9.3.2） \\
\hline
ピアノ系の統合・動作確認 & 5/28 & 6/2 & 完了 & 100\％ & ピアノ子機 \texttt{hack-ko2.ino}／Processing & 親機 → IR → ピアノ子機 → Processing の経路で譜面送信・発音を実機確認．ドラムは未テスト（付録 9.3.3） \\
\hline
\end{longtable}

%=====================================================
\section{4. 評価事項}

計画書 v1.2「1.5 テスト項目（修正版）」のうち，今回の実装（210 ピアノのプログラム，240 ドラムのプログラム，通信経路）に対応する項目を整理する．

\subsection*{4.1 対応する有効指標 MOE}

\begin{longtable}{|p{1.2cm}|p{4cm}|p{8cm}|}
\hline
\rowcolor[gray]{0.9}
No. & MOE & 評価基準 \\
\hline
MOE3 & 楽器音らしい音色 & 合成音がピアノとして聴衆に認識される \\
\hline
\end{longtable}

\subsection*{4.2 対応する性能指標 MOP（計画書 v1.2）}

\begin{longtable}{|p{1.2cm}|p{3.6cm}|p{9.0cm}|}
\hline
\rowcolor[gray]{0.9}
No. & テスト名 & 内容 \\
\hline
MOP1 & 音色単体テスト（各楽器） & Processing 単体で各楽器音を出力・確認 \\
\hline
MOP2 & シリアル通信テスト & Arduino--Processing 間通信の正常動作確認 \\
\hline
MOP3 & 赤外線通信テスト & 親機→子機の通信成功率 95\％以上 \\
\hline
MOP5 & フェードアウトテスト\textsuperscript{新規} & 音切り替え時にノイズが生じないこと \\
\hline
MOP6 & 倍速演奏テスト\textsuperscript{新規} & スイッチによるテンポ切り替え動作確認 \\
\hline
\end{longtable}

\subsection*{4.3 今週の自己評価}

\begin{longtable}{|p{1.2cm}|p{3.5cm}|p{2cm}|p{8cm}|}
\hline
\rowcolor[gray]{0.9}
No. & 項目 & 達成状況 & 備考 \\
\hline
MOE3 & 楽器音らしい音色 & OK & 倍音 [1.0，0.6，0.35，0.2，0.1，0.05] + ADSR で発音．聴感でピアノとして認識可 \\
\hline
MOP1 & 音色単体テスト & OK & Processing 単体で C4〜A4 の発音を確認 \\
\hline
MOP2 & シリアル通信テスト & OK & 921600bps でピアノ子機 → Processing の MIDI 受信．取りこぼしなく発音（ドラム経路は未テスト） \\
\hline
MOP3 & 赤外線通信テスト & 暫定OK & 親機 → ピアノ子機経由で譜面を送信し発音を確認．ドラムのプログラムは未テスト．成功率の定量計測も未実施 \\
\hline
MOP5 & フェードアウト & 未対応 & \texttt{noteOff()} が即時 \texttt{unpatch} のため切り替えノイズの懸念．第8週で対応 \\
\hline
MOP6 & 倍速演奏 & 未対応 & 親機側スイッチとの連携が未実装 \\
\hline
\end{longtable}

\subsection*{4.4 今後評価予定の項目（参考）}

\begin{itemize}
    \item MOP3：赤外線通信成功率 95\％以上の定量測定（取りこぼし時の挙動含む）
    \item MOP4：音色認識テスト（聴衆過半数が各楽器と識別）
    \item MOP5：音切り替え時のフェードアウト実装後に再評価
    \item MOP6：親機スイッチ連携後のテンポ切り替え確認
\end{itemize}

%=====================================================
\section{5. 課題管理}

\begin{longtable}{|p{2.6cm}|p{2.8cm}|p{1cm}|p{1cm}|p{3.0cm}|p{1.6cm}|p{1.6cm}|}
\hline
\rowcolor[gray]{0.9}
課題 & 発生原因 & 影響度 & 優先度 & 対策 & 期限 & 状態 \\
\hline
音切り替え時のノイズ & \texttt{noteOff()} が即時 \texttt{unpatch}（MOP5 未対応） & 高 & 高 & \texttt{Line} で 0.03 秒フェードアウト後に切断 & 6/16 & 未対応 \\
\hline
赤外線通信成功率が未計測 & MOP3（95\％目標）の定量試験未実施 & 中 & 中 & 一定回数の送受信でログ集計し成功率を算出 & 6/16 & 未対応 \\
\hline
取りこぼし時の無音・ズレ & ドラム子機は自前タイマを持たず親機任せ & 中 & 中 & 親機の \texttt{IR\_GAP\_MS} 調整と再送方針の検討 & 6/23 & 未対応 \\
\hline
\end{longtable}

\vspace{0.5em}
※影響度＝プロジェクトへのインパクト\\
※優先度＝対応の緊急性

%=====================================================
\section{6. AI利用状況と検証}

\subsection*{6.1 利用内容}

\begin{itemize}
    \item 利用目的：Processing 側のシリアル受信→MIDI→周波数変換→発音の実装，およびドラム子機 \texttt{.ino} の NEC 受信処理の実装
    \item 入力（プロンプト要旨）：
    \begin{itemize}
        \item 親機が NEC で「\texttt{address}＝子機 ID，\texttt{command}＝MIDI 番号」を送る前提を提示し，Processing 側の受信・発音コードを依頼
        \item 既存試作（第5週の \texttt{hackathon.pde}，音名文字列ベース）を MIDI 番号ベースへ移行する方針を指示
        \item IRremote v4.x（NEC）でドラム子機が \texttt{address==2} のみ拾う最小実装を依頼
    \end{itemize}
    \item 出力概要：
    \begin{itemize}
        \item Processing 側 \texttt{hack\_ko2.pde}（MIDI 受信→\texttt{PianoInstrument} 発音）
        \item ドラム子機 \texttt{hack-drum.ino}（NEC 受信→シリアル中継）
    \end{itemize}
\end{itemize}

\subsection*{6.2 検証方法}

\begin{itemize}
    \item 検証手法：
    \begin{itemize}
        \item Processing 単体で MIDI 番号を流し込み，C4〜A4 が正しい音高で鳴るか聴感確認
        \item ピアノ子機（\texttt{hack-ko2.ino}）と親機（\texttt{hack-oya2.ino}）を接続し，IR 経由で届く譜面で発音されるか実機確認（ドラムのプログラムは未テスト）
    \end{itemize}
    \item 参照資料：
    \begin{itemize}
        \item 計画書・設計書 v1.2（2026/5/29 版）「1.2 PBS」「2.2 詳細設計」
        \item Minim ライブラリ公式ドキュメント（\texttt{ddf.minim.ugens}）
        \item IRremote ライブラリ（v4.x，NEC プロトコル）
    \end{itemize}
    \item 検証手順：
    \begin{enumerate}
        \item 親機・各子機・Processing を接続して起動
        \item 親機の \texttt{melody[]} と \texttt{drumPattern[]} に沿って IR 送出
        \item ピアノ子機が MIDI 番号をシリアルへ中継し，Processing が発音するかを聴感と波形で確認（ドラムのプログラムは同手順で第8週にテスト予定）
    \end{enumerate}
\end{itemize}

\subsection*{6.3 ファクトチェック結果}

\begin{itemize}
    \item 一致点：
    \begin{itemize}
        \item MIDI→周波数変換 $f=440\cdot2^{(m-69)/12}$ が標準的な平均律の式と一致
        \item シリアル速度 921600bps が設計書「2.1 基本設計の修正」の確定値と一致
        \item NEC の \texttt{address}／\texttt{command} の対応（子機 ID／MIDI 番号）が親機実装と整合
    \end{itemize}
    \item 不一致・疑義：
    \begin{itemize}
        \item 設計書 2.1 が要求する「\texttt{noteOff} で 0.03 秒フェードアウト後に切断」が未実装（現状は即時 \texttt{unpatch}）．MOP5 の課題として残存
    \end{itemize}
    \item 最終判断：採用（今週の到達点としては妥当）
    \item 根拠：
    \begin{itemize}
        \item 親機 → IR → ピアノ子機 → Processing の経路で譜面送信・発音まで確認でき，ピアノ系の目的を満たしているため（ドラムのプログラムは第8週にテスト予定）
    \end{itemize}
\end{itemize}

%=====================================================
\section{7. 次回計画}

\begin{itemize}
    \item 目標（第8週）：
    \begin{itemize}
        \item ドラムのプログラム（\texttt{hack-drum.ino}，240）の実機テスト：親機 → ドラム子機 → Processing で，久家のドラム音色（\texttt{.pde}）と結合して受信・発音を確認
        \item 120（通信用関数・ピアノ）の正式実装と整理（受信パースの堅牢化）
        \item MOP5：\texttt{noteOff()} に \texttt{Line} フェードアウト（0.03 秒）を実装し，切り替えノイズを除去
        \item MOP3：赤外線通信成功率の定量測定（95\％目標に対する実測）
    \end{itemize}
    \item 達成基準：
    \begin{itemize}
        \item 音切り替え時にノイズが聞こえないこと
        \item 通信成功率を数値で提示できること
        \item 他楽器子機と合わせた同期演奏が破綻しないこと
    \end{itemize}
    \item 期限：
    \begin{itemize}
        \item 6/16（第8週末）
    \end{itemize}
\end{itemize}

%=====================================================
\section{8. その他}

\begin{itemize}
    \item ドラムはプログラム（櫛濱・240，\texttt{hack-drum.ino}）と音色（久家・170/171）で担当が分かれるため，\texttt{drumPattern[]} の最終仕様と \texttt{.pde} 側の発音をすり合わせ予定
    \item 親機担当（渡邊）と \texttt{IR\_GAP\_MS} や倍速演奏（MOP6）スイッチの連携仕様を確認する
\end{itemize}

%=====================================================
\section{9. 付録}

\subsection*{9.1 添付資料}

\begin{itemize}
    \item ソースコード：\texttt{hack\_ko2.pde}（ピアノ Processing プログラム），\texttt{hack-drum.ino}（ドラムのプログラム・240）
    \item ドラム音色の \texttt{.pde}（170/171）は久家担当のため本ログには含めない
    \item 参考：\texttt{hack-oya2.ino}（親機・統合用），\texttt{hack-ko2.ino}（ピアノ子機 Arduino）
    \item リポジトリ：\url{https://github.com/k-kushihama/hackathon-code/}
\end{itemize}

\subsection*{9.2 添付範囲}

\begin{itemize}
    \item 210（ピアノのプログラム）に対応する Processing 側ソースの現時点スナップショット
    \item 担当したドラムのプログラム（\texttt{hack-drum.ino}）のソース，および通信経路の構成図・回路図
\end{itemize}

\subsection*{9.3 ソースコード・図}

\subsubsection*{9.3.1 ピアノ Processing プログラム（\texttt{hack\_ko2.pde}）}

子機 Arduino からシリアルで届く MIDI 番号を受信し，$f=440\cdot2^{(m-69)/12}$ で周波数へ変換して \texttt{PianoInstrument} で発音する．音色は第5週に作成した倍音構造（\texttt{WavetableGenerator.gen10} の振幅配列 [1.0, 0.6, 0.35, 0.2, 0.1, 0.05]）と \texttt{Line} による ADSR をそのまま用いる．

\begin{lstlisting}[language=Java]
import processing.serial.*;
Serial port;
import ddf.minim.*;
import ddf.minim.ugens.*;
Waveform currentWaveform;
Minim minim;
AudioOutput out;
int currentNote = 0;

void setup() {
  size(512, 200);
  minim = new Minim(this);
  out = minim.getLineOut();
  out.setTempo(120);
  currentWaveform = WavetableGenerator.gen10(4096, new float[] { 1.0f, 0.6f, 0.35f, 0.2f, 0.1f, 0.05f });
  port = new Serial(this, "/dev/cu.usbmodem34B7DA64C6002", 921600);
  port.clear();
  port.bufferUntil('\n');
}

void draw() {
  background(0);
  stroke(255);
  for (int i = 0; i < out.bufferSize() - 1; i++) {
    line(i, 50 + out.left.get(i)*50, i+1, 50 + out.left.get(i+1)*50);
  }
  fill(255);
  text("note: " + currentNote, 10, 180);
}

void serialEvent(Serial p) {
  String inString = p.readStringUntil('\n');
  if (inString != null) {
    inString = trim(inString);
    if (inString.length() == 0) return;
    int midi = int(inString);
    currentNote = midi;
    if (midi > 0) {
      playNote(midi);
    }
  }
}

void playNote(int midi) {
  float freq = (float)(440.0 * Math.pow(2.0, (midi - 69) / 12.0));
  out.playNote(0, 0.4f, new PianoInstrument(freq, 0.6f, currentWaveform));
}

class PianoInstrument implements Instrument {
  Oscil wave;
  Line ampEnv;
  Oscil vib;
  Summer freqSum;
  Constant baseFreq;
  float maxAmp;

  PianoInstrument(float frequency, float maxAmp, Waveform wf) {
    this.maxAmp = maxAmp;
    freqSum = new Summer();
    baseFreq = new Constant(frequency);
    vib = new Oscil(5, 0, Waves.SINE);
    baseFreq.patch(freqSum);
    vib.patch(freqSum);
    wave = new Oscil(frequency, 0, wf);
    freqSum.patch(wave.frequency);
    ampEnv = new Line();
    ampEnv.patch(wave.amplitude);
  }

  void noteOn(float duration) {
    ampEnv.activate(duration, this.maxAmp, 0);
    wave.patch(out);
  }

  void noteOff() {
    wave.unpatch(out);
  }
}
\end{lstlisting}

\subsubsection*{9.3.2 ドラムのプログラム（\texttt{hack-drum.ino}・GC の 240）}

ドラムは作業を分担し，本プログラム（子機側の NEC 受信→シリアル中継）を櫛濱が，ドラム音色の \texttt{.pde}（170/171）を久家が担当する．赤外線受光モジュールで NEC フレームを受信し，自分宛（\texttt{CHILD\_ID = 2}）のものだけを拾って \texttt{command}（MIDI 番号）をシリアルへ中継する．楽譜を持たずステートレスで，テンポはすべて親機任せのため可変 BPM に自動追従する設計だが，第6週に対面できず\textbf{実機テストは未実施}．

\begin{lstlisting}[language=C++]
#define DECODE_NEC
#include <IRremote.hpp>

#define CHILD_ID 2

const int PIN_IR_RECV = 2;
const int LED_INDICATOR = 13;

void setup() {
  Serial.begin(921600);
  IrReceiver.begin(PIN_IR_RECV);
  pinMode(LED_INDICATOR, OUTPUT);
}

void loop() {
  if (IrReceiver.decode()) {
    if (IrReceiver.decodedIRData.protocol == NEC) {
      uint16_t address = IrReceiver.decodedIRData.address;
      if (address == CHILD_ID) {
        uint8_t note = (uint8_t)IrReceiver.decodedIRData.command;
        Serial.println(note);
        digitalWrite(LED_INDICATOR, HIGH);
        delay(15);
        digitalWrite(LED_INDICATOR, LOW);
      }
    }
    IrReceiver.resume();
  }
}
\end{lstlisting}

\subsubsection*{9.3.3 システム統合：データフロー}

図~\ref{fig:dataflow} に，親機から発音までのデータフローを示す．親機（\texttt{hack-oya2.ino}）が \texttt{melody[]}／\texttt{drumPattern[]} に沿って NEC フレーム（\texttt{address}＝子機 ID，\texttt{command}＝MIDI 番号）を送出し，各子機は自分宛のフレームだけを拾って MIDI 番号をシリアルへ中継，Processing が受信して発音する．ピアノ子機は \texttt{CANON\_OFFSET\_TICKS} 分だけ遅れて開始することで輪唱（カノン）を実現する．本ログ時点では，このうち\textbf{ピアノ系（親機 → ピアノ子機 → Processing）の譜面送信・発音までを実機で確認}しており，ドラム経路（\texttt{hack-drum.ino} ＋ 久家のドラム音色 \texttt{.pde}）は未テストである．

\begin{figure}[H]
\centering
\begin{tikzpicture}[
  node distance=0.6cm and 1.0cm,
  block/.style={rectangle, draw, rounded corners=2pt, minimum width=2.7cm, minimum height=1.1cm, align=center, font=\small},
  io/.style={rectangle, draw, fill=gray!10, minimum width=2.3cm, minimum height=1.0cm, align=center, font=\small},
  arrow/.style={-Stealth, thick}
]
  \node[block] (oya) {親機 Arduino\\\texttt{hack-oya2.ino}\\(NEC 送出)};
  \node[block, right=of oya] (ko) {子機 Arduino\\\texttt{hack-ko2/hack-drum}\\(NEC 受信→中継)};
  \node[block, right=of ko] (proc) {Processing\\\texttt{hack\_ko2.pde}\\(MIDI→発音)};
  \node[io, right=of proc] (spk) {スピーカー};

  \draw[arrow] (oya) -- node[above, font=\scriptsize]{IR / NEC} node[below, font=\scriptsize]{addr=ID, cmd=MIDI} (ko);
  \draw[arrow] (ko) -- node[above, font=\scriptsize]{USB Serial} node[below, font=\scriptsize]{921600bps} (proc);
  \draw[arrow] (proc) -- node[above, font=\scriptsize]{audio} (spk);
\end{tikzpicture}
\caption{システム統合のデータフロー（親機 → IR/NEC → 子機 → USB シリアル → Processing → 発音）}
\label{fig:dataflow}
\end{figure}

\subsubsection*{9.3.4 ドラム子機 Arduino の回路図}

図~\ref{fig:circuit} にドラム子機 Arduino の接続を示す．赤外線受光モジュール CHQ1838D を 1 つ用い，親機からの NEC フレームを受信する．D2 に受光信号，D13 に動作確認用 LED を接続し，USB を Processing 実行用 PC へつなぐ．

\begin{figure}[H]
\centering
\begin{tikzpicture}[
  every node/.style={font=\small},
  pin/.style={circle, draw, fill=black, inner sep=1.2pt},
  module/.style={rectangle, draw, minimum width=2.4cm, minimum height=1.2cm, align=center},
  arduino/.style={rectangle, draw, very thick, rounded corners=4pt, minimum width=3.8cm, minimum height=5.2cm, align=center},
  wire/.style={thick}
]
  \node[arduino] (ard) at (0,0) {Arduino\\UNO R4\\(子機: ドラム)};

  \node[pin, label=right:{\tiny D2}] (d2) at (1.9,1.2) {};
  \node[pin, label=right:{\tiny D13(LED)}] (d13) at (1.9,0.0) {};
  \node[pin, label=right:{\tiny 5V}] (vcc) at (1.9,-1.0) {};
  \node[pin, label=right:{\tiny GND}] (gnd) at (1.9,-2.0) {};
  \node[pin, label=left:{\tiny USB}] (usb) at (-1.9,0) {};

  \node[module] (ir1) at (5.4,1.2) {IR 受光\\CHQ1838D\\(NEC 受信)};

  \node[pin] (ir1out) at (4.2,1.2) {};
  \node[pin] (ir1vcc) at (4.2,1.6) {};
  \node[pin] (ir1gnd) at (4.2,0.8) {};

  \draw[wire] (d2) -- (ir1out);

  \draw[wire] (vcc) -- ++(0.7,0) coordinate (vbus) -- (vbus|-ir1vcc) -- (ir1vcc);
  \draw[wire] (gnd) -- ++(1.1,0) coordinate (gbus) -- (gbus|-ir1gnd) -- (ir1gnd);

  \node[rectangle, draw, dashed, minimum width=2cm, minimum height=1cm, align=center] (pc) at (-4.5,0) {PC\\(Processing)};
  \draw[wire, -Stealth] (usb) -- node[above, font=\scriptsize]{USB / Serial} (pc.east);
\end{tikzpicture}
\caption{ドラム子機 Arduino の回路図．D2 に IR 受光モジュール，D13 に動作確認 LED を接続し，USB 経由で Processing 実行用 PC と通信する．}
\label{fig:circuit}
\end{figure}

なお，赤外線受光モジュールは内部にプルアップを持つため，Arduino 側は \texttt{INPUT\_PULLUP} 相当で十分．電源は Arduino の 5V／GND を共用する．

\end{document}
